\chapter{Derivatives: One and Several Variables}

A position logger on a bicycle writes one point every second. To obtain the velocity, we
divide the difference in distance by the difference in time. Two readings one second apart
and four meters apart give four meters per second, even though that is only an average
velocity. If the separation is one hundredth of a second, the difference in distance is a few
centimeters while the instrument noise is several meters; the quotient jumps to hundreds of
meters per second, even with a negative sign. The same division, two different answers.

The derivative resolves that tension by moving the question to the limit, where noise does not
enter the count. But the limit cannot be measured: what is measurable is only the quotient at
a finite separation, and it carries two errors at once, truncation because the separation is
not zero and rounding because the machine stores only so many digits. The derivative is the
exact mathematical object; the quotient is its approximation, with an error that can be
computed and controlled.

The second difficulty appears when there is more than one input. A credit scoring model
receives twelve inputs, and shifting one input gives one number for the rate of change.
Shifting two inputs at once is not the same as adding those two numbers, unless the function
behaves well; later we shall see an example that fails, and that function is still
continuous. The correct answer is not a list of twelve numbers, but a single linear map on the
shift vector: that is the true derivative, and its matrix is the Jacobian. The continuity of
Chapter 11 does not say how fast something changes.

\begin{untukanak}
Have you ever stood on a curved roof? If you step forward, you climb. If you step sideways,
you can stay at the same height. And yet you are standing at the same point, on the same roof.
So the question ``how steep is this place'' is not complete before we name the direction of
the step.

There is another problem. If you look at only two points on the roof, and then divide the
difference in height by the difference in distance, you get an average slope, not the slope at
a single point. To get the slope exactly under your feet, the two points must be brought
closer and closer. The closer they are the more exact it is, but if they are too close your
measurement numbers fall apart because the measuring instrument is not perfect.

In this chapter we shall learn why the slope at a single point needs a direction, why the
steepest direction always exists, and why a computing tool whose step is too tight is
precisely wrong.
\end{untukanak}

\section{The derivative as a limit and its two meanings}
\tingkat{basic}

The instantaneous rate of change is the number approached by the quotient of two differences
as the difference in the input shrinks. The following definition merely writes that sentence
down carefully, because the quotient itself is undefined when its denominator is zero.

\kamus{derivative}{turunan}{The slope of the line that hugs the graph most tightly at a single
point, which is also the instantaneous rate of change at that point.}

\begin{definisi}[The derivative at a single point]
Suppose $f$ is defined on an open interval containing $c$. The derivative of $f$ at $c$ is
\[
  f'(c)=\lim_{h\to0}\frac{f(c+h)-f(c)}{h},
\]
provided that limit exists and is finite. In this case $f$ is called
\emph{differentiable} at $c$. Other notation used in this book: $Df(c)$ and
$\dfrac{df}{dx}(c)$.\index{derivative}
\end{definisi}

That quotient is defined for every $h\neq0$, so there is no division by zero; what may fail
is the limit. The failure is real: for $f(x)=|x|$ at $c=0$ the quotient is
$1$ for $h>0$ and $-1$ for $h<0$, so the two-sided limit does not exist.

Physically, the derivative is the instantaneous rate of change: $s'(t)$ is velocity when
$s(t)$ is position, $Q'(t)$ is current when $Q(t)$ is the charge that flows, and $C'(q)$ is
marginal cost when $C(q)$ is the cost of producing $q$ units.

Geometrically, the quotient is the slope of the secant from $(c,f(c))$ to
$(c+h,f(c+h))$. As $h\to0$, that secant rotates about the first point toward a single
position; the line in the limiting position is called the tangent line, and its slope is
$f'(c)$. The equation of the tangent line at $c$ is
\[
  y=f(c)+f'(c)(x-c).
\]
The tangent line is not ``the line that touches the graph exactly once'', for with
$f(x)=x^{3}$ at $c=0$ the tangent line $y=0$ crosses the graph at the origin. What is correct:
its difference from the graph shrinks faster than the distance does.

That rotation can be followed numerically. For $f(x)=x^{2}$ at $c=2$, the secant with $h=1$
has slope $5$, with $h=0.5$ the slope is $4.5$, and the tangent line has slope $4$.
Figure~\ref{fig:sekan-singgung} draws all three: from $5$ to $4.5$, settling at $4$.

\begin{figure}[htbp]
\centering
\begin{tikzpicture}
\begin{axis}[
  width=\linewidth, height=6.4cm,
  axis lines=left,
  xlabel={$x$}, ylabel={$y$},
  xmin=1.55, xmax=3.25, ymin=2.2, ymax=10.2,
  xtick={2,3}, ytick={4,9},
  tick label style={font=\footnotesize},
  label style={font=\footnotesize},
  legend style={font=\scriptsize, draw=none, fill=none,
                at={(0.03,0.97)}, anchor=north west},
  clip=false,
]
\addplot[ink, thick, domain=1.6:3.15, samples=90] {x^2};
\addlegendentry{$f(x)=x^{2}$}
\addplot[accent, thick, domain=1.95:3.18] {5*(x-2)+4};
\addlegendentry{secant $h=1$, slope $5$}
\addplot[layerteal, thick, dashed, domain=1.95:3.0] {4.5*(x-2)+4};
\addlegendentry{secant $h=0.5$, slope $4.5$}
\addplot[layerindigo, thick, domain=1.68:3.05] {4*(x-2)+4};
\addlegendentry{tangent line, slope $4$}
\addplot[only marks, mark=*, mark size=1.5pt, ink]
  coordinates {(2,4) (2.5,6.25) (3,9)};
\node[font=\scriptsize, anchor=north west, text=ink] at (axis cs:2.04,3.75) {$(2,4)$};
\end{axis}
\end{tikzpicture}
\caption{Secants rotating into the tangent line of $f(x)=x^{2}$ at $(2,4)$.
Their slopes drop from $5$ (for $h=1$) to $4.5$ (for $h=0.5$) and settle at $4$
as $h\to0$.}
\label{fig:sekan-singgung}
\end{figure}

\begin{bacakeras}
Two voices take turns; the answer gets shorter every time.

--- How fast is it moving? --- From which point to which point?

--- From the first second to the second second. --- Three meters.

--- From the first second to second one and a half. --- Two and a half meters.

--- From the first second to second one and one tenth. --- Two and one fifth.

--- From the first second to second one and one hundredth. --- Almost two.

--- From the first second to the first second. --- I cannot. The denominator is zero.

Stop there. What changes is not the answer, but the distance between the two points being
compared. The distance shrinks, the secant rotates, the number settles: a value that is never
reached yet always approached, that is the derivative.

Repeat together: the distance shrinks, the line rotates, the number settles.
\end{bacakeras}

Let us compute one derivative from the definition. For $f(x)=x^{2}$ at $c=2$,
\[
  \frac{f(2+h)-f(2)}{h}=\frac{(4+4h+h^{2})-4}{h}=\frac{h(4+h)}{h}=4+h .
\]
Since $h\neq0$, that division is legitimate, and the limit is $4$. For $h=0.01$ the quotient
is $4.01$; the difference $0.01$ is not an arithmetic mistake but the order-$h$ term that
disappears only in the limit. For $f(x)=1/x$ at $c=2$,
\[
  \frac{1/(2+h)-1/2}{h}
  =\frac{2-(2+h)}{2h(2+h)}=\frac{-1}{2(2+h)}\longrightarrow-\frac14 .
\]
In both computations the algebra removes the factor $h$ from the denominator, so that the
form $0/0$ becomes a function continuous at $h=0$.

That quotient form is important, but it does not survive when the variables become many. The
form that survives is the linear approximation.

\begin{definisi}[Derivative as the coefficient of a linear approximation]
A function $f$ is differentiable at $c$ with derivative $L$ if
\[
  f(c+h)=f(c)+Lh+r(h),
  \qquad\text{with}\qquad
  \lim_{h\to0}\frac{r(h)}{h}=0 .
\]
\end{definisi}

The remainder $r(h)$ is smaller than $h$; that is the content of the definition. For one
variable this form feels excessive, but only this form survives up to the Jacobian section.

\begin{proposisi}
The two definitions of the derivative above are equivalent, and the number $L$ in the second
definition is unique.
\end{proposisi}

\begin{proof}
If the limit of the quotient equals $L$, set $r(h)=f(c+h)-f(c)-Lh$. Then $r(h)/h$ is exactly
the quotient minus $L$, hence tends to $0$. Conversely, if $r$ with that property exists, the
quotient equals $L+r(h)/h$, which tends to $L$.

For uniqueness, suppose $L_{1}$ and $L_{2}$ both satisfy it. Then
$(L_{1}-L_{2})h=r_{2}(h)-r_{1}(h)$, and after dividing by $h$ we obtain
$L_{1}-L_{2}=r_{2}(h)/h-r_{1}(h)/h$. The right-hand side tends to $0$, whereas the left-hand
side does not depend on $h$. Hence $L_{1}=L_{2}$.
\end{proof}

The first property to hold on to: differentiability is stronger than continuity.

\begin{proposisi}
If $f$ is differentiable at $c$, then $f$ is continuous at $c$.
\end{proposisi}

\begin{proof}
For $h\neq0$ we have
\[
  f(c+h)-f(c)=h\cdot\frac{f(c+h)-f(c)}{h}.
\]
The second factor tends to $f'(c)$ and the first factor tends to $0$, so the product tends to
$0$. In other words $\lim_{h\to0}f(c+h)=f(c)$.
\end{proof}

The converse does not hold. The function $|x|$ is continuous at $0$ but not differentiable
there, because the two one-sided quotients have opposite signs. The function $x\sin(1/x)$,
equal to $0$ at $x=0$, is also continuous but not differentiable there: its quotient is
$\sin(1/x)$, which oscillates without approaching a single value.

\begin{tantangan}
Suppose $f(0)=0$ and, for $x\neq0$,
\[
  f(x)=|x|\bigl(2+\sin(1/x)\bigr).
\]
Show that $f$ is continuous at $0$ but has no derivative there. The two properties do not
contradict each other: continuity demands that the values of the function approach $f(0)$,
the derivative demands that the difference quotient approach a single number.

Hint 1: for continuity, one inequality that traps $|f(x)|$ between $0$ and something tending
to $0$ is enough.

Hint 2: for the second, compute $f(h)/h$ separately for $h>0$ and $h<0$. For $h>0$ the
quotient is $2+\sin(1/h)$, which oscillates between $1$ and $3$; for $h<0$ it oscillates
between $-3$ and $-1$. The two sets of values do not overlap.

If you get stuck, go back to the definition of the derivative at the beginning of this section
and compute the quotient for $h>0$ and then $h<0$ for the function $|x|$.
\end{tantangan}

The following computation rules free us from the definition every time.

\begin{teorema}[Basic rules]
Suppose $f$ and $g$ are differentiable at $c$. Then
\begin{enumerate}
  \item $(f+g)'(c)=f'(c)+g'(c)$ and $(af)'(c)=a f'(c)$ for $a\in\mathbb{R}$;
  \item $(fg)'(c)=f'(c)g(c)+f(c)g'(c)$;
  \item $\left(\dfrac{f}{g}\right)'(c)=\dfrac{g(c)f'(c)-f(c)g'(c)}{g(c)^{2}}$ if
        $g(c)\neq0$;
  \item if $n$ is an integer, the derivative of $x^{n}$ at $x\neq0$ is $nx^{n-1}$.
\end{enumerate}
\end{teorema}

\begin{proof}
Part (1) is immediate from the limit of a sum. For (2), write
\[
  f(x)g(x)-f(c)g(c)=f(x)\bigl[g(x)-g(c)\bigr]+g(c)\bigl[f(x)-f(c)\bigr].
\]
After dividing by $x-c$ and letting $x\to c$, the factor $f(x)$ tends to $f(c)$ because $f$ is
continuous at $c$ (the previous proposition), and both quotients tend to their derivatives.

For (3), it suffices to show that $1/g$ is differentiable with derivative $-g'(c)/g(c)^{2}$,
then apply (2) to $f\cdot(1/g)$. The quotient for $1/g$ is
\[
  \frac{1/g(x)-1/g(c)}{x-c}
  =-\frac{1}{g(x)g(c)}\cdot\frac{g(x)-g(c)}{x-c},
\]
which tends to $-g'(c)/g(c)^{2}$ because $g$ is continuous and $g(c)\neq0$. Part (4) follows
from (2) by induction on $n\ge0$, then from (3) for $n<0$.
\end{proof}
\section{The chain rule}
\tingkat{intermediate}

If one quantity depends on a second quantity, and the second quantity depends on time, then
the rate of change of the first quantity is the product of the two rates. That sentence sounds
like everyday arithmetic, and that is indeed the content of the chain rule.

Figure~\ref{fig:roda-gigi} uses two gears: the inner wheel turns three times per turn of the
middle wheel, and the middle wheel twice per turn of the outer wheel. One turn of the outer
wheel turns the inner wheel $3\times2=6$ times; adding the two ratios gives $5$, a number that
appears in no turn at all.

\begin{figure}[htbp]
\centering
\begin{tikzpicture}[scale=0.95]
% inner wheel, ratio 3
\draw[ink, thick] (0,0) circle (1.05);
\foreach \a in {0,20,...,340} { \draw[ink, thick] (\a:1.05) -- (\a:1.17); }
\fill[ink] (0,0) circle (1.5pt);
\draw[->, layerindigo, thick] (0,0) ++(150:0.55) arc (150:70:0.55);
% outer wheel, ratio 2
\begin{scope}[xshift=2.9cm]
  \draw[ink, thick] (0,0) circle (1.8);
  \foreach \a in {10,30,...,350} { \draw[ink, thick] (\a:1.8) -- (\a:1.95); }
  \fill[ink] (0,0) circle (1.5pt);
  \draw[->, layerbrown, thick] (0,0) ++(-20:0.75) arc (-20:-100:0.75);
\end{scope}
\node[font=\footnotesize, text=ink, align=center] at (0,-2.35) {inner wheel:\\3 turns};
\node[font=\footnotesize, text=ink, align=center] at (2.9,-2.35) {outer wheel:\\2 turns};
\node[font=\footnotesize, text=layerplum] at (1.45,2.15) {$3\times2=6$};
\end{tikzpicture}
\caption{The chain rule as two gears. The inner wheel turns $3$ times for one turn of the
outer wheel, and the outer wheel $2$ times for one turn of its output; the total ratio is
$3\times2=6$, not $3+2=5$.}
\label{fig:roda-gigi}
\end{figure}

\kamus{chain rule}{aturan rantai}{The rate at the end of the chain equals the product of the
rates at every link.}

\begin{teorema}[Chain rule]
Suppose $g$ is differentiable at $c$ and $f$ is differentiable at $g(c)$. Then $f\circ g$
is differentiable at $c$ and
\[
  (f\circ g)'(c)=f'\bigl(g(c)\bigr)\,g'(c).
\]
\end{teorema}

\begin{proof}
Write $y=g(c)$. Since $f$ is differentiable at $y$, by the linear approximation definition
there is a function $v$ with $v(s)\to0$ as $s\to y$ and
\[
  f(s)-f(y)=(s-y)\bigl[f'(y)+v(s)\bigr].
\]
In the same way, since $g$ is differentiable at $c$, there is $u$ with $u(t)\to0$ as
$t\to c$ and
\[
  g(t)-g(c)=(t-c)\bigl[g'(c)+u(t)\bigr].
\]
Substitute $s=g(t)$ in the first equation, then use the second:
\[
  f\bigl(g(t)\bigr)-f\bigl(g(c)\bigr)
  =(t-c)\bigl[g'(c)+u(t)\bigr]\bigl[f'(y)+v(g(t))\bigr].
\]
For $t\neq c$, divide by $t-c$. Since $g$ is continuous at $c$, $g(t)\to y$, so that
$v(g(t))\to0$, and the quotient tends to $g'(c)f'(y)$.
\end{proof}

The naive proof that divides by $g(t)-g(c)$ fails when $g(t)=g(c)$ for $t$ near $c$ but
$t\neq c$; a locally constant function, for instance, makes that denominator zero at every
point. The buffer $v$ above avoids that division altogether.

Short example: for $h(x)=(x^{2}+1)^{5}$, take $g(x)=x^{2}+1$ and $f(u)=u^{5}$. Then
$h'(x)=5(x^{2}+1)^{4}\cdot2x$ and $h'(1)=5\cdot16\cdot2=160$; direct expansion is possible,
but unpleasant.

\begin{resep}
\begin{enumerate}
  \item \textbf{Separate the two layers.} Point out which is inside and which is outside. For
        $k(x)=\sin(x^{3}+2x)$: the inner layer is $g(x)=x^{3}+2x$, the outer layer
        $f(u)=\sin u$.
  \item \textbf{Differentiate the outer layer, leave its content intact.} $f'(u)=\cos u$
        becomes $\cos(x^{3}+2x)$; do not simplify the content yet.
  \item \textbf{Multiply by the derivative of the inner layer.} $g'(x)=3x^{2}+2$, so that
        $k'(x)=(3x^{2}+2)\cos(x^{3}+2x)$.
  \item \textbf{Test one point.} At $x=0$ the formula gives $2$, and the central difference
        with $h=10^{-3}$ gives $1.9999997$.
\end{enumerate}
Warning about order: the outer layer is evaluated at $g(x)$, not at $x$. Swapping the place of
evaluation gives $\cos(x)\,(3x^{2}+2)$, whose value at $x=1$ is $2.7015$, whereas the true
value is $-4.9499$.
\end{resep}

Rate example: the area of a circle $A=\pi r^{2}$ with $\dot r=0.5$ cm per second. At
$r=10$ cm,
\[
  \frac{dA}{dt}=2\pi r\frac{dr}{dt}=2\pi\cdot10\cdot0.5=10\pi\approx31.4159
  \ \text{cm}^{2}\ \text{per second}.
\]
The number $10\pi$ is exact; at $r=20$ cm the same rate gives $20\pi$: the answer depends on
$r$.

\section{The mean value theorem and its consequences}
\tingkat{intermediate}

Between any two points on a smooth graph there is always a point whose slope equals the slope
of the secant joining those two points. That statement looks like geometric fantasy, yet it is
the bridge between the derivative and the function.

The key to its proof is the special case in which the two ends have equal value.

\begin{teorema}[Rolle]
If $f$ is continuous on $[a,b]$, differentiable on $(a,b)$, and $f(a)=f(b)$, then there is
$c\in(a,b)$ with $f'(c)=0$.\index{mean value theorem}
\end{teorema}

\begin{proof}
If $f$ is constant, every $c\in(a,b)$ works. If not, $f$ attains a maximum or a minimum that
differs from the value at the ends; since $f$ is continuous on a closed bounded interval,
that value is attained at some point, and that point cannot be an endpoint, so $c\in(a,b)$
(extreme value theorem, Chapter 11). At an interior extreme point, the quotient
$(f(c+h)-f(c))/h$ is $\le0$ for small $h>0$ and $\ge0$ for small $h<0$. The two inequalities
together force its two-sided limit, namely $f'(c)$, to be $0$.
\end{proof}

\begin{teorema}[Mean value]
If $f$ is continuous on $[a,b]$ and differentiable on $(a,b)$, then there is $c\in(a,b)$
with
\[
  f'(c)=\frac{f(b)-f(a)}{b-a}.
\]
\end{teorema}

\begin{proof}
Set $s=(f(b)-f(a))/(b-a)$ and $g(x)=f(x)-f(a)-s(x-a)$. Then $g$ is continuous on $[a,b]$,
differentiable on $(a,b)$, and $g(a)=g(b)=0$. By Rolle's theorem there is $c\in(a,b)$ with
$g'(c)=0$, that is, $f'(c)=s$.
\end{proof}

Numerical example: $f(x)=x^{3}$ on $[0,3]$ has secant slope $27/3=9$, and $3c^{2}=9$ gives
$c=\sqrt{3}\approx1.7321$; there is no guarantee that $c$ sits in the middle.

\begin{tangan}
Prepare a thin 30 cm board, two books, a ruler, and a pencil. Prop up the left end of the
board to a height of 4 cm, then press the middle slightly so that it curves upward.

Hold the ruler upright at 5 cm and at 25 cm from the raised end; the difference in height
divided by 20 cm gives one slope number. Repeat with 5 cm and 10 cm: the number changes,
because what is being measured is the average slope over different intervals.

Now lay the ruler flat at the 15 cm point and rotate it until only one point touches. Lift the
other end slowly: the ruler rotates about that point, exactly like the secant in
Figure~\ref{fig:sekan-singgung}. The slope of the ruler that touches a single point is the
slope at that point, and that is where the mean value theorem promises that such a point
exists.
\end{tangan}

\begin{akibat}[Monotonicity test]
Suppose $f$ is continuous on $[a,b]$ and differentiable on $(a,b)$. If $f'(x)\ge0$ for all
$x\in(a,b)$, then $f$ is nondecreasing on $[a,b]$. If $f'(x)>0$ for all $x$, then $f$ is
strictly increasing. If $f'(x)=0$ for all $x$, then $f$ is constant.
\end{akibat}

\begin{proof}
Take $a\le x<y\le b$. The mean value theorem on $[x,y]$ gives $c\in(x,y)$ with
$f(y)-f(x)=f'(c)(y-x)$. Since $y-x>0$, the sign of the difference in values follows the sign
of $f'(c)$. If $f'$ is zero everywhere, the difference is zero for every pair, so $f$ is
constant.
\end{proof}

The direction of the logic deserves attention: this corollary turns information about the
derivative into information about the function. The converse does not hold: a strictly
increasing function may have zero derivative at some points, such as $x^{3}$ at $0$. Nor does
the conclusion run backward from a single point: $f'(c)=0$ does not mean that $c$ is an
extreme point.

\begin{akibat}[Lipschitz continuity]
If $f$ is differentiable on an interval $I$ and $|f'(x)|\le M$ for all $x\in I$, then
$|f(x)-f(y)|\le M|x-y|$ for all $x,y\in I$.
\end{akibat}

\begin{proof}
For $x\neq y$, the mean value theorem on the interval between them gives $c$
with $f(x)-f(y)=f'(c)(x-y)$. Take the absolute value of both sides and use $|f'(c)|\le M$.
\end{proof}

Therefore a function with bounded derivative is automatically uniformly continuous (take
$\delta=\varepsilon/M$), but the converse does not hold: $\sqrt{x}$ is uniformly continuous on
$[0,1]$ with an unbounded derivative near zero.

\section{Taylor polynomials and their remainder}
\tingkat{advanced}

The tangent line is a polynomial of degree one that clings to the graph at a single point. If
we want it to cling longer, we add terms of degree two, three, and so on. The question is not
whether the additions improve the approximation, but how large the remainder is.

\begin{definisi}[Taylor polynomial]
Suppose $f$ has derivatives up to order $n$ at $a$. The Taylor polynomial of order $n$ of $f$
at $a$ is
\[
  P_{n}(x)=\sum_{k=0}^{n}\frac{f^{(k)}(a)}{k!}(x-a)^{k}.
\]
\end{definisi}

The requirement: $P_{n}$ and $f$ agree in value and in their first $n$ derivatives at $a$;
since $n+1$ coefficients determine one polynomial of degree $n$, $P_{n}$ is unique.

\begin{teorema}[Taylor with Lagrange remainder]
Suppose $f$ has derivatives up to order $n-1$ that are continuous on $[a,b]$ and a derivative
of order $n$ on $(a,b)$. For every $\beta\in(a,b]$ there is $\xi\in(a,\beta)$ with
\[
  f(\beta)=P_{n-1}(\beta)+\frac{f^{(n)}(\xi)}{n!}(\beta-a)^{n}.
\]
\end{teorema}

\begin{proof}
Define $M$ by $f(\beta)=P_{n-1}(\beta)+M(\beta-a)^{n}$, then set
\[
  g(t)=f(t)-P_{n-1}(t)-M(t-a)^{n},\qquad a\le t\le b .
\]
$P_{n-1}$ is built from the derivatives of $f$ at $a$, so $g(a)=\cdots=g^{(n-1)}(a)=0$, and
the choice of $M$ gives $g(\beta)=0$. Rolle gives $x_{1}\in(a,\beta)$ with $g'(x_{1})=0$;
since $g'(a)=0$, Rolle applied to $g'$ on $[a,x_{1}]$ gives $x_{2}$ with $g''(x_{2})=0$.
Repeat through $n$ steps: there is $x_{n}\in(a,\beta)$ with $g^{(n)}(x_{n})=0$. The derivative
of order $n$ of $P_{n-1}$ is zero and that of $(t-a)^{n}$ is $n!$, so
$0=f^{(n)}(x_{n})-n!M$, that is, $n!M=f^{(n)}(\xi)$ with $\xi=x_{n}$.
\end{proof}

For $n=1$ this theorem is exactly the mean value theorem; what is new is the ability to
compute errors of higher order.

\begin{contoh}
Approximate $e$ by $P_{3}(1)$ from $f(x)=e^{x}$ at $a=0$, then compare with the remainder
bound.
\end{contoh}

Here $f^{(k)}(0)=1$ for every $k$, so that
\[
  P_{3}(1)=1+1+\frac{1}{2}+\frac{1}{6}=\frac{8}{3}=2.6666\ldots
\]
The true value is $e=2.71828\ldots$, so the error is $0.0516152$. The Lagrange remainder
gives $|e-P_{3}(1)|=e^{\xi}/24$ for some $\xi\in(0,1)$, and since $e^{\xi}\le e<3$, the bound
is $3/24=0.125$. That bound is correct, but loose by a factor of about two and a half. Using
order five, $P_{5}(1)=2.71666\ldots$ with error $0.0016152$, whereas the bound is
$3/720=0.0041667$.

\begin{untukanak}
Imagine that you are guessing your friend's route from a few clues only. First clue: where he
is now. Your guess is a single point, and that point does not move. Second clue: which way he
is stepping. Your guess is a straight line, already closer, but a real route usually curves.
Third clue: how fast the direction changes. Your guess is a piece of a parabola, which clings
longer.

The more clues, the longer your guess clings, but never exactly. There is always a remainder.
That remainder can be bounded without knowing the true route, only from how large the change
is at the next level. In this chapter we shall learn why a straight line is the guess of order
one, why a parabola is of order two, and why a remainder bound is more useful than an answer
that claims to be exact.
\end{untukanak}
\section{Partial derivatives, the gradient, and directional derivatives}
\tingkat{intermediate}

When there is more than one input, ``rate of change'' is not necessarily a single number. It
depends on the direction in which we step. The simplest way to obtain a single number is to
freeze every coordinate except one, and then use the definition we already have.

\begin{definisi}[Partial derivative]
Suppose $f\colon U\to\mathbb{R}$ is defined on an open set
$U\subseteq\mathbb{R}^{n}$ and $a\in U$. The partial derivative of $f$ with respect to $x_{j}$
at $a$ is
\[
  \frac{\partial f}{\partial x_{j}}(a)
  =\lim_{h\to0}\frac{f(a+h\,e_{j})-f(a)}{h},
\]
with $e_{j}$ the unit vector in the direction of $x_{j}$, provided the limit exists.
\index{partial derivative}
\end{definisi}

Geometrically, the partial derivative with respect to $x$ is the slope of the slice of the
graph by the plane $y$ constant, hence an ordinary derivative of a function of one variable.

\kamus{gradient}{gradien}{An arrow pointing in the steepest direction on a hill, with length
equal to the steepness in that direction.}

\begin{definisi}[Gradient]
If all partial derivatives of $f$ at $a$ exist, the vector
\[
  \nabla f(a)=\left(\frac{\partial f}{\partial x_{1}}(a),\dots,
  \frac{\partial f}{\partial x_{n}}(a)\right)
\]
is called the gradient of $f$ at $a$.\index{gradient}
\end{definisi}

For $f(x,y)=x^{2}y^{3}$ we have $\partial f/\partial x=2xy^{3}$ and
$\partial f/\partial y=3x^{2}y^{2}$, so that $\nabla f(2,-1)=(-4,12)$.

\begin{definisi}[Directional derivative]
For a unit vector $u\in\mathbb{R}^{n}$, the directional derivative of $f$ at $a$ along $u$ is
\[
  D_{u}f(a)=\lim_{t\to0}\frac{f(a+tu)-f(a)}{t},
\]
provided the limit exists.\index{directional derivative}
\end{definisi}

The partial derivatives are the special case $u=e_{j}$. For other directions $D_{u}f$ cannot
be predicted from the list of partials, except in one well-behaved situation.

\begin{proposisi}
If $f$ is differentiable at $a$ (in the sense of the linear approximation of the next section),
then for every unit vector $u$ we have $D_{u}f(a)=\nabla f(a)\cdot u$.
\end{proposisi}

\begin{proof}
With $v=tu$ in the linear approximation definition,
$f(a+tu)=f(a)+t\,\nabla f(a)\cdot u+r(tu)$, so that
\[
  \frac{f(a+tu)-f(a)}{t}=\nabla f(a)\cdot u+\frac{r(tu)}{t}.
\]
Since $\|r(tu)\|/\|tu\|\to0$ and $\|tu\|=|t|$, the last term tends to $0$.
\end{proof}

\begin{proposisi}[Direction of steepest ascent]
If $\nabla f(a)\neq0$, then among all unit vectors $u$ the value $D_{u}f(a)$ is largest at
$u=\nabla f(a)/\|\nabla f(a)\|$, and the largest value is $\|\nabla f(a)\|$.
\end{proposisi}

\begin{proof}
By the Cauchy--Schwarz inequality (Chapter 10),
\[
  |D_{u}f(a)|=|\nabla f(a)\cdot u|\le\|\nabla f(a)\|\,\|u\|=\|\nabla f(a)\|.
\]
Equality holds exactly when $u$ is parallel to $\nabla f(a)$. Since $u$ is a unit vector, only
two signs are possible along that direction, and the positive sign gives the largest value.
\end{proof}

For $f(x,y)=x^{2}y^{3}$ at $(2,-1)$, the length of the gradient is $\sqrt{160}\approx12.6491$,
so the direction of steepest ascent is $(-4,12)/\sqrt{160}\approx(-0.3162,\,0.9487)$ with
rate $12.6491$ per unit of distance. The $x$ direction gives $-4$, the $y$ direction gives
$12$; both are smaller.

If a curve $t\mapsto\gamma(t)$ stays on the level set $f=\text{constant}$, then
$f(\gamma(t))$ is constant and its derivative is zero; the chain rule gives
$\nabla f(\gamma(t))\cdot\gamma'(t)=0$. Hence the gradient is perpendicular to the level
curve, and the direction of steepest ascent is perpendicular to the contour lines on a map.

For example $f(x,y)=x^{2}+2y^{2}$ at $(2,-1)$, a point with $f=6$, so it sits on the level
ellipse $x^{2}+2y^{2}=6$. Its gradient is $(2x,4y)=(4,-4)$, of length
$4\sqrt{2}\approx5.6569$, whereas the tangent direction of the ellipse there is $(1,1)$ and
$(4,-4)\cdot(1,1)=0$. Figure~\ref{fig:kontur-gradien} draws all three: the level ellipse, the
gradient arrow, and the tangent line perpendicular to it. The gradient arrow does not run
along the contour; it pierces the contour in the direction in which the value of $f$
increases.

\begin{figure}[htbp]
\centering
\begin{tikzpicture}[scale=0.92]
\draw[->, layergrey] (-3.0,0) -- (3.4,0) node[right, font=\scriptsize] {$x$};
\draw[->, layergrey] (0,-2.45) -- (0,2.5) node[above, font=\scriptsize] {$y$};
\draw[ink] (0,0) ellipse [x radius=1.414, y radius=1.000];
\draw[ink] (0,0) ellipse [x radius=2.000, y radius=1.414];
\draw[ink, thick] (0,0) ellipse [x radius=2.449, y radius=1.732];
\node[font=\scriptsize, anchor=west, text=ink] at (0.28,-0.98) {$c=2$};
\node[font=\scriptsize, anchor=west, text=ink] at (0.28,-1.40) {$c=4$};
\node[font=\scriptsize, anchor=west, text=ink] at (0.28,-1.98) {$c=6$};
\draw[accent, thick, dashed] (1.30,-1.70) -- (2.72,-0.28);
\draw[->, very thick, layerindigo] (2,-1) -- (2.707,-1.707);
\fill[ink] (2,-1) circle (1.7pt);
\node[font=\scriptsize, anchor=north west, text=ink] at (2.02,-1.05) {$(2,-1)$};
\node[font=\scriptsize, anchor=west, text=layerindigo] at (2.72,-1.86) {$\nabla f=(4,-4)$};
\node[font=\scriptsize, anchor=east, text=accent] at (1.32,-1.76) {direction $(1,1)$};
\end{tikzpicture}
\caption{Level curves $x^{2}+2y^{2}=c$ for $c=2,4,6$ and the gradient at $(2,-1)$. The
gradient arrow $(4,-4)$ is perpendicular to the tangent direction $(1,1)$ of the level curve
$c=6$, because $(4,-4)\cdot(1,1)=0$.}
\label{fig:kontur-gradien}
\end{figure}

\section{Differentiability, the Jacobian, and the multivariable chain rule}
\tingkat{advanced}

The true derivative is not a list of partials, but a single linear map that imitates the
function near a point; the list of partials can mislead.

\begin{definisi}[Differentiability]
Suppose $U\subseteq\mathbb{R}^{n}$ is open and $F\colon U\to\mathbb{R}^{m}$. $F$ is
differentiable at $a\in U$ if there is a linear map $L\colon\mathbb{R}^{n}\to\mathbb{R}^{m}$
with
\[
  F(a+v)=F(a)+Lv+r(v),
  \qquad
  \lim_{v\to0}\frac{\|r(v)\|}{\|v\|}=0 .
\]
The map $L$ is called the derivative of $F$ at $a$ and is written $F'(a)$.
\index{differentiability}
\end{definisi}

\begin{proposisi}
The map $L$ in that definition is unique if it exists.
\end{proposisi}

\begin{proof}
If $L_{1}$ and $L_{2}$ both work, then $(L_{1}-L_{2})v=r_{2}(v)-r_{1}(v)$. Take
$v=tw$ with $w$ fixed and $t>0$, divide by $t$, then let $t\to0$: this gives
$(L_{1}-L_{2})w=0$. Since this holds for every $w$, the two maps are equal.
\end{proof}

\begin{teorema}
If $F$ is differentiable at $a$, then all its partial derivatives at $a$ exist, and the matrix
of $L$ with respect to the standard bases is
\[
  J_{F}(a)=\begin{pmatrix}
    \dfrac{\partial F_{1}}{\partial x_{1}}(a) & \cdots &
    \dfrac{\partial F_{1}}{\partial x_{n}}(a)\\
    \vdots & \ddots & \vdots\\
    \dfrac{\partial F_{m}}{\partial x_{1}}(a) & \cdots &
    \dfrac{\partial F_{m}}{\partial x_{n}}(a)
  \end{pmatrix}.
\]
\end{teorema}

\begin{proof}
Take $v=te_{j}$ with $t\neq0$. The definition gives
$F(a+te_{j})-F(a)=tLe_{j}+r(te_{j})$, and after dividing by $t$ and letting $t\to0$, the $j$-th
column of $L$ appears as the limit of a quotient, namely the vector of partial derivatives
with respect to $x_{j}$.
\end{proof}

\kamus{Jacobian matrix}{Jacobian}{A table of all the slopes of a function of several variables:
each row is one output, each column is one input.}

That $m\times n$ matrix is called the Jacobian matrix of $F$ at $a$.\index{jacobian} For $m=1$
it becomes a single row, namely the transpose of the gradient. Example:
$F(x,y)=(x^{2}y,\;x+y^{2})$ at $(1,2)$ gives
\[
  J_{F}(1,2)=\begin{pmatrix}2xy & x^{2}\\ 1 & 2y\end{pmatrix}_{(1,2)}
  =\begin{pmatrix}4&1\\ 1&4\end{pmatrix},
\]
so that $F(1.01,\,1.98)\approx(2,5)+J_{F}(1,2)(0.01,-0.02)=(2.02,\,4.93)$,
whereas the exact value is $(2.019798,\,4.9304)$; the approximation is off by
$2\times10^{-4}$ in the first component.

Existence of the partials does not guarantee differentiability; the next two examples show two
levels of failure.

\begin{contoh}[Partials exist, function discontinuous]
Let $f(x,y)=\dfrac{xy}{x^{2}+y^{2}}$ for $(x,y)\neq(0,0)$ and $f(0,0)=0$. At the origin
$f(h,0)=0$, so both partial derivatives vanish there. But along the line $y=kx$ the value of
the function is $k/(1+k^{2})$: for $k=1$ the value is $1/2$, for $k=0$ the value is $0$. The
limit does not exist, and $f$ is not continuous at the origin.
\end{contoh}

\begin{contoh}[Continuous, partials exist, not differentiable]
Let $g(x,y)=\dfrac{xy}{\sqrt{x^{2}+y^{2}}}$ for $(x,y)\neq(0,0)$ and $g(0,0)=0$.
Since $2|xy|\le x^{2}+y^{2}$, we have $|g(x,y)|\le\frac12\sqrt{x^{2}+y^{2}}$, so $g$ is
continuous at the origin, and both of its partials vanish there. If $g$ were differentiable,
its linear map would have to be zero, so the remainder $r(v)=g(v)$ would have to satisfy
$\|r(v)\|/\|v\|\to0$. Take $v=t(1,1)/\sqrt{2}$ with $t>0$: then $g(v)=(t^{2}/2)/t=t/2$, so
$g(v)/\|v\|=1/2$ for every $t$. That ratio tends to $1/2$, not to $0$; hence $g$ is not
differentiable at the origin.
\end{contoh}

The consequence matters: in the second example $\nabla g(0,0)=0$ yet the diagonal directional
derivative does not exist, because its quotient is $\pm1/2$ depending on the sign of $t$. Thus
$D_{u}f=\nabla f\cdot u$ requires differentiability.

A condition that is easy to check and sufficient for differentiability is continuity of the
partials.

\begin{teorema}
If the partial derivatives of $F$ exist and are continuous near $a$, then $F$ is
differentiable at $a$.
\end{teorema}

\begin{proof}
It suffices to prove this for $m=1$ and $n=2$; the other cases follow component by component
with increasing indices. Write $a=(p,q)$ and $v=(h,k)$. Then
\[
  f(p+h,q+k)-f(p,q)
  =\bigl[f(p+h,q+k)-f(p,q+k)\bigr]+\bigl[f(p,q+k)-f(p,q)\bigr].
\]
By the mean value theorem in one variable, the first term equals
$h\,\partial_{x}f(\xi,q+k)$ for some $\xi$ between $p$ and $p+h$, and the second term equals
$k\,\partial_{y}f(p,\eta)$ for some $\eta$ between $q$ and $q+k$. Therefore
\[
  f(a+v)-f(a)-\nabla f(a)\cdot v
  =h\bigl[\partial_{x}f(\xi,q+k)-\partial_{x}f(a)\bigr]
  +k\bigl[\partial_{y}f(p,\eta)-\partial_{y}f(a)\bigr].
\]
Divide by $\|v\|=\sqrt{h^{2}+k^{2}}$: since $|h|/\|v\|\le1$ and $|k|/\|v\|\le1$, and both
differences in brackets tend to $0$ (this is where continuity of the partials is used), the
right-hand side divided by $\|v\|$ tends to $0$.
\end{proof}

Now composition.

\begin{teorema}[Multivariable chain rule]
Suppose $H\colon\mathbb{R}^{n}\to\mathbb{R}^{k}$ is differentiable at $a$ and
$G\colon\mathbb{R}^{k}\to\mathbb{R}^{m}$ is differentiable at $H(a)$. Then $G\circ H$
is differentiable at $a$ and
\[
  (G\circ H)'(a)=G'(H(a))\,H'(a),
\]
that is, the composition of two linear maps; in matrices,
$J_{G\circ H}(a)=J_{G}(H(a))\,J_{H}(a)$.
\end{teorema}

\begin{proof}
Write $A=H'(a)$, $B=G'(y_{0})$ with $y_{0}=H(a)$, and
\[
  H(a+v)=H(a)+Av+r(v),\qquad G(y_{0}+w)=G(y_{0})+Bw+q(w),
\]
with $\|r(v)\|\le\varepsilon(v)\|v\|$ and $\|q(w)\|\le\eta(w)\|w\|$, where
$\varepsilon(v)\to0$ and $\eta(w)\to0$. For $w=H(a+v)-H(a)=Av+r(v)$ we have
$\|w\|\le(\|A\|+\varepsilon(v))\|v\|$, so $w\to0$ as $v\to0$. Then
\[
  G\bigl(H(a+v)\bigr)-G(H(a))-BAv
  =B\bigl(H(a+v)-H(a)-Av\bigr)+q(w)
  =B\,r(v)+q(w).
\]
Divide by $\|v\|$:
\[
  \frac{\|B\,r(v)+q(w)\|}{\|v\|}
  \le\|B\|\,\varepsilon(v)+\bigl(\|A\|+\varepsilon(v)\bigr)\eta(w),
\]
and the right-hand side tends to $0$. Hence $BA$ satisfies the definition of the derivative of
$G\circ H$ at $a$.
\end{proof}

The matrix form explains why the order of multiplication may not be reversed: $J_{H}$ is
$k\times n$ and $J_{G}$ is $m\times k$, so only one order can be multiplied. For scalar
functions the result is $\frac{d}{dt}f(\gamma(t))=\nabla f(\gamma(t))\cdot
\gamma'(t)$.

Finally, second derivatives of a function of several variables are stored in a single matrix.

\kamus{Hessian matrix}{Hessian}{A table of curvature: whether the surface bends upward,
downward, or in both directions at once.}

\begin{definisi}[Hessian matrix]
If all second partial derivatives of $f\colon U\to\mathbb{R}$ exist at $a$, the matrix
\[
  H_{f}(a)=\left(\frac{\partial^{2}f}{\partial x_{i}\,\partial x_{j}}(a)\right)_{i,j=1}^{n}
\]
is called the Hessian matrix of $f$ at $a$.\index{hessian}
\end{definisi}

For a function whose second partials are continuous, $H_{f}(a)$ is symmetric (Clairaut's
theorem, proved in Stewart in the section on partial derivatives). Without continuity, that
conclusion can fail.

The Hessian is the derivative of the gradient. If the second partials are continuous near a
critical point $a$ ($\nabla f(a)=0$), then
\[
  f(a+v)=f(a)+\tfrac12\,v^{\top}H_{f}(a)\,v+o\bigl(\|v\|^{2}\bigr),
\]
so that the sign of the quadratic form $v\mapsto v^{\top}Hv$ determines the type of the
critical point.

\begin{proposisi}[Second derivative test]
Suppose $f$ has continuous second partial derivatives near a critical point $a$
($\nabla f(a)=0$).
\begin{enumerate}
  \item If $H_{f}(a)$ is positive definite, $a$ is a strict local minimum.
  \item If $H_{f}(a)$ is negative definite, $a$ is a strict local maximum.
  \item If $H_{f}(a)$ has eigenvalues of different signs, $a$ is a saddle point.
  \item If $H_{f}(a)$ is semidefinite but singular, this test gives no conclusion.
\end{enumerate}
\end{proposisi}

The first three parts follow from the quadratic expansion above: along a direction with
$v^{\top}Hv>0$ the function increases, along a direction with $v^{\top}Hv<0$ it decreases. The
complete proof is in Stewart in the section on maximum and minimum values. The fourth part is
genuinely empty: $x^{4}+y^{4}$ and $x^{4}-y^{4}$ both have $H=0$ at the origin, but the first
is a minimum and the second a saddle. For two variables, $AC-B^{2}>0$ with $A>0$ means
positive definite, $AC-B^{2}<0$ means indefinite.
\section{Worked examples}
\tingkat{intermediate}

The following three computations use the same numbers as the code in the lab section.

First, finite differences for $f(x)=\sin x$ at $x_{0}=1$; the exact value is $\cos 1$, whose
digits appear in the code output in the lab section. The two formulas:
\[
  \text{forward: }\frac{f(x_{0}+h)-f(x_{0})}{h},\qquad
  \text{central: }\frac{f(x_{0}+h)-f(x_{0}-h)}{2h}.
\]

\begin{center}
\begin{tabular}{rrrrr}
\hline
$h$ & forward difference & relative error & central difference & relative error\\
\hline
$10^{-1}$ & $0.497363752535$ & $7.947\times10^{-2}$ & $0.539402252170$ & $1.666\times10^{-3}$\\
$10^{-2}$ & $0.536085981012$ & $7.804\times10^{-3}$ & $0.540293300875$ & $1.667\times10^{-5}$\\
$10^{-4}$ & $0.540260231419$ & $7.787\times10^{-5}$ & $0.540302304968$ & $1.667\times10^{-9}$\\
$10^{-6}$ & $0.540301885121$ & $7.787\times10^{-7}$ & $0.540302305896$ & $5.130\times10^{-11}$\\
$10^{-8}$ & $0.540302302898$ & $5.497\times10^{-9}$ & $0.540302308449$ & $4.777\times10^{-9}$\\
$10^{-10}$ & $0.540302247387$ & $1.082\times10^{-7}$ & $0.540302247387$ & $1.082\times10^{-7}$\\
\hline
\end{tabular}
\end{center}

The pattern of the errors is no accident. For the forward difference, the Taylor expansion
gives
\[
  \frac{f(x_{0}+h)-f(x_{0})}{h}=f'(x_{0})+\frac{h}{2}f''(x_{0})+\mathcal{O}(h^{2}),
\]
so its relative error is of order $h$ and drops by a factor of ten each time $h$ is divided by
ten: $7.947\times10^{-2}$, $7.804\times10^{-3}$, $7.787\times10^{-5}$. The numbers can be
matched: $\frac{h}{2}f''(1)=-\frac{1}{2}\cdot10^{-2}\sin 1=-4.207\times10^{-3}$, and
$4.207\times10^{-3}/0.5403=7.787\times10^{-3}$. The central difference gives
\[
  \frac{f(x_{0}+h)-f(x_{0}-h)}{2h}=f'(x_{0})+\frac{h^{2}}{6}f'''(x_{0})+\mathcal{O}(h^{4}),
\]
so its error is of order $h^{2}$ and drops by a factor of a hundred: $1.666\times10^{-3}$,
$1.667\times10^{-5}$, $1.667\times10^{-9}$. Here $\frac{10^{-4}}{6}\cos 1
=9.005\times10^{-6}$, and $9.005\times10^{-6}/0.5403$ again gives $1.667\times10^{-5}$.

At $h=10^{-8}$ and $10^{-10}$ the central difference deteriorates again: the truncation error
is already smaller than the rounding error, which grows like $\varepsilon/h$, with
$\varepsilon\approx2.2\times10^{-16}$. The best entry in the table, $h=10^{-6}$ with error
$5\times10^{-11}$, has an order consistent with $\varepsilon^{1/3}\approx6\times10^{-6}$.

Second, the gradient of $f(x,y)=x^{2}y^{3}$ at $(2,-1)$ with analytic value $(-4,12)$. The
central difference gives an error of $3.333\times10^{-5}$ in the $y$ component for
$h=10^{-2}$, dropping to $3.333\times10^{-9}$ for $h=10^{-4}$: order $h^{2}$. The $x$
component is exact to machine error, because $f$ is quadratic in $x$.

Third, the chain rule on a composition of different dimensions. Let
\[
  H(x,y)=\bigl(x^{2}+y,\;xy,\;y^{3}\bigr),\qquad
  G(u,v,w)=\bigl(uv+w,\;u^{2}-vw\bigr),
\]
and consider $F=G\circ H$ at $a=(1,2)$. Here $H(a)=(3,2,8)$, and
\[
  J_{H}(1,2)=\begin{pmatrix}2&1\\ 2&1\\ 0&12\end{pmatrix}_{3\times2},
  \qquad
  J_{G}(3,2,8)=\begin{pmatrix}2&3&1\\ 6&-8&-2\end{pmatrix}_{2\times3}.
\]
Their product is $2\times2$:
\[
  J_{G}(3,2,8)\,J_{H}(1,2)=
  \begin{pmatrix}2&3&1\\ 6&-8&-2\end{pmatrix}
  \begin{pmatrix}2&1\\ 2&1\\ 0&12\end{pmatrix}
  =\begin{pmatrix}10&17\\ -4&-26\end{pmatrix}.
\]
\kertas{For the composition above, compute $J_{F}$ at $(2,1)$ in two ways: the product
$J_{G}(H(2,1))J_{H}(2,1)$ and the direct derivative of the explicit form. Then check the
direction $v=(1,3)$ by comparing $J_{F}(2,1)v$ against $\bigl(F(a+sv)-F(a)\bigr)/s$ for
$s=10^{-2}$ and $s=10^{-4}$. Why does the error drop only by a factor of $10$, not $100$?}

\section{Computational lab}
\tingkat{advanced}

The following code is run as is with numpy 2.5.2 on Python 3: the first part repeats the three
computations above, then one experiment that cannot be done by hand, finite differences under
measurement noise.

\begin{lab}
\begin{verbatim}
import numpy as np

# 1. finite differences for f(x) = sin x at x0 = 1
x0 = 1.0
exact = float(np.cos(x0))
print("cos(1) =", exact)
for h in [1e-1, 1e-2, 1e-4, 1e-6, 1e-8, 1e-10]:
    fwd = (np.sin(x0 + h) - np.sin(x0)) / h
    cen = (np.sin(x0 + h) - np.sin(x0 - h)) / (2*h)
    em = abs(fwd - exact) / abs(exact)
    ep = abs(cen - exact) / abs(exact)
    print(f"{h:8.0e} fwd={fwd:.12f} {em:.3e} "
          f"cen={cen:.12f} {ep:.3e}")
\end{verbatim}
Actual output:
\begin{verbatim}
cos(1) = 0.5403023058681398
   1e-01 fwd=0.497363752535 7.947e-02 cen=0.539402252170 1.666e-03
   1e-02 fwd=0.536085981012 7.804e-03 cen=0.540293300875 1.667e-05
   1e-04 fwd=0.540260231419 7.787e-05 cen=0.540302304968 1.667e-09
   1e-06 fwd=0.540301885121 7.787e-07 cen=0.540302305896 5.130e-11
   1e-08 fwd=0.540302302898 5.497e-09 cen=0.540302308449 4.777e-09
   1e-10 fwd=0.540302247387 1.082e-07 cen=0.540302247387 1.082e-07
\end{verbatim}
\end{lab}

The last line deserves a second reading: at $h=10^{-10}$ both formulas give identical digits,
because they are both computing the same rounding error, not the derivative.

\begin{lab}
\begin{verbatim}
def f(p):
    return p[0]**2 * p[1]**3

def grad_fd(fun, a, h):
    g = np.zeros_like(a)
    for j in range(a.size):
        e = np.zeros_like(a)
        e[j] = h
        g[j] = (fun(a + e) - fun(a - e)) / (2*h)
    return g

a = np.array([2.0, -1.0])
g_an = np.array([2*a[0]*a[1]**3, 3*a[0]**2*a[1]**2])
print("gradient =", g_an)
for h in [1e-2, 1e-4, 1e-6]:
    g = grad_fd(f, a, h)
    print(f"{h:8.0e} {g[0]:.12f} {g[1]:.12f} "
          f"{abs(g[0]+4)/4:.3e} {abs(g[1]-12)/12:.3e}")
\end{verbatim}
Actual output:
\begin{verbatim}
gradient = [-4. 12.]
   1e-02 -4.000000000000 12.000400000000 1.577e-14 3.333e-05
   1e-04 -4.000000000004 12.000000039998 1.000e-12 3.333e-09
   1e-06 -4.000000000115 11.999999999679 2.876e-11 2.676e-11
\end{verbatim}
\end{lab}

\begin{lab}
\begin{verbatim}
def hh(p):
    x, y = p
    return np.array([x*x + y, x*y, y**3])

def gg(q):
    u, v, w = q
    return np.array([u*v + w, u*u - v*w])

def F(p):
    return gg(hh(p))

def jac(fun, p, step):
    n = p.size
    J = np.zeros((fun(p).size, n))
    for j in range(n):
        e = np.zeros(n)
        e[j] = step
        J[:, j] = (fun(p + e) - fun(p - e)) / (2*step)
    return J

a = np.array([1.0, 2.0])
q = hh(a)
Jh = np.array([[2*a[0], 1.0], [a[1], a[0]], [0.0, 3*a[1]**2]])
Jg = np.array([[q[1], q[0], 1.0], [2*q[0], -q[2], -q[1]]])
J_chain = Jg @ Jh
J_num = jac(F, a, 1e-6)
print("h(1,2) =", q)
print("Jg @ Jh =\n", J_chain)
print("max difference =", f"{np.max(np.abs(J_num - J_chain)):.3e}")

v = np.array([0.7, -1.3])
d_chain = J_chain @ v
d_num = (F(a + 1e-6*v) - F(a - 1e-6*v)) / 2e-6
print("J v =", d_chain)
print("finite difference =", d_num)
print("relative error =",
      f"{np.max(np.abs(d_num - d_chain)/np.abs(d_chain)):.3e}")
\end{verbatim}
Actual output:
\begin{verbatim}
h(1,2) = [3. 2. 8.]
Jg @ Jh =
 [[ 10.  17.]
 [ -4. -26.]]
max difference = 1.858e-09
J v = [-15.1  31. ]
finite difference = [-15.1  31. ]
relative error = 9.860e-11
\end{verbatim}
\end{lab}

The difference $1.858\times10^{-9}$ between the chain-rule Jacobian and the finite-difference
Jacobian is the truncation error of the central difference, not an algebraic mistake; the
directional check gives $9.86\times10^{-11}$, because both use the same point.

\begin{lab}
\begin{verbatim}
rng = np.random.default_rng(7)
sigma = 1e-3
for h in [1e-1, 1e-2, 1e-3, 1e-4, 1e-5]:
    err = []
    for _ in range(2000):
        s1 = np.sin(x0 + h) + sigma*rng.standard_normal()
        s0 = np.sin(x0) + sigma*rng.standard_normal()
        err.append(abs((s1 - s0)/h - exact))
    print(f"{h:8.0e} mean error = {np.mean(err):.3e}")
\end{verbatim}
Actual output:
\begin{verbatim}
   1e-01 mean error = 4.203e-02
   1e-02 mean error = 1.112e-01
   1e-03 mean error = 1.160e+00
   1e-04 mean error = 1.117e+01
   1e-05 mean error = 1.088e+02
\end{verbatim}
\end{lab}

The last experiment is the heart of this chapter in numbers. With noise
$\sigma=10^{-3}$, the quotient carries two errors: a truncation bias of $\approx0.42h$ and
noise that amplifies like $\sqrt{2}\sigma/h$. At $h=10^{-1}$ the bias dominates ($0.042$),
at $h=10^{-4}$ the noise dominates ($11.2$), and the two balance when
$0.42h=\sqrt{2}\cdot10^{-3}/h$, that is $h\approx0.058$: with an instrument that bad, the
best step is about six hundredths, not one hundred-millionth.

That is the honest answer to the opening question. The derivative is a limit, but its best
approximation is not the closest possible limit. The best is a compromise between two errors
that move in opposite directions, and that compromise can be computed before a single number
is measured.
\begin{untukdosen}
The order that works: the quotient definition, then the linear approximation form, then the
computation rules. That form feels excessive for one variable, but only it survives up to the
Jacobian; if it is introduced later, students memorize the chain rule as matrix multiplication
without knowing why the order is what it is.

Two typical misconceptions. First, ``the partial derivatives exist means the function is
smooth''; the example $xy/\sqrt{x^{2}+y^{2}}$ on the board breaks it, provided the facts are
computed together: partials zero, continuous, yet the diagonal direction gives $1/2$. Second,
the direction of multiplication in the chain rule: students write $J_{H}J_{G}$ because the
order in which the functions are written is the reverse of the order in which they are
applied; the way to check is the size of the matrices.

Blackboard activity: ask the class to compute both partials at the origin for the two
counterexamples, then to decide whether that function has a tangent plane. Sign of readiness
to go further: students explain why $D_{u}f=\nabla f\cdot u$ fails in the second example, and
use the word ``remainder'', not ``formula''.
\end{untukdosen}

\section{Summary}
\tingkat{basic}

\begin{ringkasan}
\begin{itemize}
  \item The derivative $f'(c)$ is the limit of a difference quotient, and also the unique
        coefficient of the linear approximation $f(c+h)=f(c)+Lh+r(h)$ with $r(h)/h\to0$.
        Differentiability guarantees continuity, not the other way around.
  \item The chain rule $(f\circ g)'=f'(g(c))g'(c)$ is proved through the linear approximation,
        and the proof survives into many dimensions: $J_{G\circ H}(a)=J_{G}(H(a))J_{H}(a)$.
  \item The mean value theorem links the sign of the derivative to the behavior of the
        function: $f'\ge0$ means nondecreasing, $|f'|\le M$ means Lipschitz; it fails for
        vector functions.
  \item Taylor with Lagrange remainder uses an unknown $\xi$, so only an upper bound for
        $|f^{(n)}|$ is used.
  \item For several variables, the derivative is a linear map with a Jacobian matrix. Partials
        may exist without continuity, and may exist together with continuity without
        differentiability.
  \item The gradient points in the direction of steepest ascent with rate $\|\nabla f\|$,
        perpendicular to the contours; a definite Hessian determines the type of a critical
        point.
  \item Finite differences are of order $h$ (forward) or $h^{2}$ (central), and the rounding
        error grows like $\varepsilon/h$. The smallest step is not the best step.
\end{itemize}
\end{ringkasan}

\section{Exercises}
\tingkat{intermediate}

\begin{latihan}
\begin{enumerate}
  \item Compute $f'(2)$ from the definition for $f(x)=x^{3}$ and $f(x)=\sqrt{x}$, until the
        factor $h$ disappears from the denominator.

  \item Prove that if $f'(c)>0$, there is $\delta>0$ such that $f(c+h)>f(c)$ for every
        $h\in(0,\delta)$, and explain why this does not mean that $f$ is increasing on some
        interval.

  \item Prove the quotient rule by differentiating $1/g$ first, then using the product rule.
        State at which step the continuity of $g$ is used.

  \item Assess this claim: \emph{if $f$ has all its partial derivatives at a point, then $f$
        is continuous at that point.} State whether it is true, false, or not yet supported by
        proof, and include a counterexample if needed.

  \item For $f(x,y)=x^{3}y-xy^{3}$, compute the gradient at $(2,1)$, determine the direction
        of steepest ascent, and compute its maximum rate. Check with a central difference
        $h=10^{-4}$ and compare the relative error with the order $h^{2}$.

  \item Prove that if $f'=0$ everywhere on an interval, then $f$ is constant. Explain why that
        proof fails on a disconnected set, and give an example.

  \item Suppose $H(x,y)=(e^{x}y,\;x+y^{2})$ and $G(u,v)=(u^{2}v,\;u-v)$. Compute
        $J_{G\circ H}$ at $(0,1)$ in two ways: the matrix product and the direct derivative.
        (Hint: $H(0,1)=(1,1)$.)

  \item Find all critical points of $f(x,y)=x^{3}+y^{3}-3xy$, then classify them with the
        Hessian; explain why the sign of the determinant is not enough when the determinant is
        zero.

  \item Estimate $\sqrt{1.04}$ with the Taylor polynomial of order two at $a=1$, then bound
        its error with the Lagrange remainder. Compare that bound with the actual error: is it tight? (Hint:
        $f''(x)=-\frac14x^{-3/2}$ and $f'''$ is positive on that interval.)
\end{enumerate}
\end{latihan}

\begin{untukpeneliti}
What is proved in this chapter are the classical statements for one and several variables: the
linear approximation and the uniqueness of its map, the chain rule, the mean value theorem and
its consequences, Taylor with Lagrange remainder, the equivalence of differentiability with the
Jacobian, and the sufficient condition of continuous partials. All of them are theorems, not
conventions, and they hold in finite-dimensional spaces over $\mathbb{R}$.

What is often popularized is the claim ``the gradient is the derivative''. Without
differentiability that claim is false, and $xy/\sqrt{x^{2}+y^{2}}$ is a counterexample that
stays continuous. The claim ``the Hessian determines the type of a critical point'' holds only
when the Hessian is nonsingular; the singular case requires higher derivatives or direct
analysis.

Primary sources. Rudin, \emph{Principles of Mathematical Analysis}, the chapters
\emph{Differentiation} and \emph{Functions of Several Variables}, contains all the statements
above. Spivak, \emph{Calculus on Manifolds}, the chapter \emph{Differentiation}, takes the same
route and is shorter. Stewart, \emph{Calculus: Early Transcendentals} ninth edition, sections
14.3 to 14.7, supplies the computational counterparts, including Clairaut's theorem.

Limits of validity: all arguments assume that the domain is open and that the derivatives used
exist. For functions that are differentiable only almost everywhere, the monotonicity and
Taylor conclusions collapse, and that is territory of measure theory which this book does not
cover. For vector functions, the mean value theorem does not hold; what replaces it is the mean
value inequality, which requires an operator norm. The best order of a finite difference
depends on the scale of the function, the noise, and the arithmetic; this chapter has only
shown the order empirically.
\end{untukpeneliti}

\peta{This chapter follows Rudin, \emph{Principles of Mathematical Analysis}, the chapters
\emph{Differentiation} and \emph{Functions of Several Variables}; its computational version is
Stewart, \emph{Calculus: Early Transcendentals} ninth edition, sections 14.3 to 14.7; its
abstract counterpart is Spivak, \emph{Calculus on Manifolds}, the chapter
\emph{Differentiation}. Continuity and the extreme value theorem are used from Chapter 11, the
Cauchy--Schwarz inequality from Chapter 10, and the Jacobian norm returns in Chapter 23 and
Chapter 28.}
